\documentclass[10pt,letterpaper]{article}

\usepackage{amsmath,amssymb,needspace}

\setlength{\textwidth}{7.35in}
\setlength{\textheight}{9.55in}
\setlength{\oddsidemargin}{-0.42in}
\setlength{\evensidemargin}{-0.42in}
\setlength{\topmargin}{-0.58in}
\setlength{\columnsep}{0.28in}
\setlength{\parindent}{0pt}
\setlength{\parskip}{2.8pt}
\raggedbottom
\sloppy

\makeatletter
\renewcommand{\tableofcontents}{%
  \section*{Contents}%
  \@starttoc{toc}%
}
\renewcommand*\l@section[2]{%
  \ifnum \c@tocdepth>\z@
    \vskip 0.42em%
    \noindent{\bfseries #1}%
    \hspace{0.55em}\nobreak\leaders\hbox{\rule[0.55ex]{0.75em}{0.35pt}}\hfill
    \nobreak\hspace{0.55em}%
    \nobreak\hb@xt@\@pnumwidth{\hfil\bfseries #2}\par
  \fi
}
\makeatother

\newcommand{\doublerule}{%
  \noindent\hbox to \columnwidth{%
    \vbox{%
      \hrule width \columnwidth height 0.36pt
      \kern 0.7pt
      \hrule width \columnwidth height 0.36pt
    }%
  }%
}

\newcommand{\lecture}[2]{%
  \par\Needspace{6\baselineskip}\vspace{0.85em}%
  \addcontentsline{toc}{section}{Lecture #1: #2}%
  \doublerule
  \par\nopagebreak\vspace{0.14em}%
  {\raggedright\Large\bfseries Lecture #1: #2\par}%
  \nopagebreak\vspace{0.14em}%
  \doublerule
  \par\vspace{0.38em}%
}
\newcommand{\topic}[1]{%
  \par\Needspace{4\baselineskip}\vspace{0.75em}%
  {\raggedright\large\bfseries #1\par}%
  \vspace{0.08em}%
  \noindent\rule{\columnwidth}{0.28pt}%
  \par\vspace{0.18em}%
}
\newcommand{\labelname}[1]{{\bfseries #1:}\quad}
\newcommand{\usecase}[1]{\par\noindent\labelname{Use}#1\par}
\newcommand{\formulas}{\par\noindent\labelname{Formula}\par\vspace{-0.45em}}
\newenvironment{sym}{%
  \par\noindent\labelname{Symbols}\vspace{-0.65em}
  \begin{list}{$\bullet$}{%
    \setlength{\leftmargin}{1.2em}
    \setlength{\itemsep}{0pt}
    \setlength{\parsep}{0pt}
    \setlength{\topsep}{1pt}
  }%
}{\end{list}}

\newcommand{\E}{\mathbb{E}}
\newcommand{\Prob}{\mathbb{P}}
\newcommand{\Var}{\operatorname{Var}}
\newcommand{\Cov}{\operatorname{Cov}}
\newcommand{\Corr}{\operatorname{Corr}}
\newcommand{\ind}{\mathbf{1}}

\begin{document}
\begin{center}
{\Large MIT 6.041SC Only Formula Notes}\\[0.25em]
{\normalsize Lectures 1--25.}
\end{center}
\vspace{0.35em}
\tableofcontents
\clearpage
\twocolumn
\small

\lecture{1}{Probability Models and Axioms}

\topic{Probability Axioms}
\usecase{Use for checking whether a probability law is valid.}
\formulas
\[
\Prob(A)\ge 0,\qquad \Prob(\Omega)=1,
\]
\[
A_i\cap A_j=\emptyset\ (i\ne j)
\quad\Longrightarrow\quad
\Prob\!\left(\bigcup_i A_i\right)=\sum_i \Prob(A_i).
\]
\begin{sym}
  \item $\Omega$ is the sample space.
  \item $A_i$ are events.
\end{sym}

\topic{Complement and Inclusion--Exclusion}
\usecase{Use when events overlap or the complement is easier.}
\formulas
\[
\Prob(A^c)=1-\Prob(A),
\]
\[
\Prob(A\cup B)=\Prob(A)+\Prob(B)-\Prob(A\cap B).
\]
\[
\Prob(\text{at least one }A_i)=1-\Prob(\text{none of the }A_i).
\]
\begin{sym}
  \item $A^c$ means event $A$ does not happen.
  \item $A\cup B$ means $A$ or $B$.
  \item $A\cap B$ means both $A$ and $B$.
\end{sym}

\topic{Useful Probability Bounds}
\usecase{Use for quick upper bounds or set inclusion.}
\formulas
\[
A\subseteq B\quad\Longrightarrow\quad \Prob(A)\le \Prob(B).
\]
\[
\Prob(A\cup B)\le \Prob(A)+\Prob(B).
\]
\[
\Prob\!\left(\bigcup_i A_i\right)\le \sum_i \Prob(A_i).
\]
\begin{sym}
  \item The last two are union bounds.
\end{sym}

\topic{Discrete Uniform Law}
\usecase{Use when all outcomes are equally likely.}
\formulas
\[
\Prob(A)=\frac{|A|}{|\Omega|}.
\]
\begin{sym}
  \item $|A|$ is the number of favorable outcomes.
  \item $|\Omega|$ is the total number of outcomes.
\end{sym}

\lecture{2}{Conditional Probability and Bayes}

\topic{Conditional Probability}
\usecase{Use when the problem says ``given $B$''.}
\formulas
\[
\Prob(A\mid B)=\frac{\Prob(A\cap B)}{\Prob(B)},\qquad \Prob(B)>0.
\]
\begin{sym}
  \item $A$ is the event wanted.
  \item $B$ is the condition/new universe.
\end{sym}

\topic{Multiplication Rule}
\usecase{Use for staged experiments or paths in a probability tree.}
\formulas
\[
\Prob(A\cap B)=\Prob(A)\Prob(B\mid A).
\]
\[
\Prob(A_1\cap\cdots\cap A_n)
=\Prob(A_1)\prod_{i=2}^n
\Prob(A_i\mid A_1\cap\cdots\cap A_{i-1}).
\]
\begin{sym}
  \item $A_i$ are events in sequence.
\end{sym}

\topic{Total Probability}
\usecase{Use when cases $A_i$ form a partition.}
\formulas
\[
\Prob(B)=\sum_i \Prob(A_i)\Prob(B\mid A_i).
\]
\begin{sym}
  \item $A_i$ are disjoint and exhaustive cases.
  \item $B$ is the event wanted.
\end{sym}

\topic{Bayes' Rule}
\usecase{Use when the question reverses evidence and cause.}
\formulas
\[
\Prob(A_i\mid B)=
\frac{\Prob(A_i)\Prob(B\mid A_i)}
{\sum_j \Prob(A_j)\Prob(B\mid A_j)}.
\]
\begin{sym}
  \item $A_i$ is one possible hidden case.
  \item $B$ is observed evidence.
  \item $\Prob(A_i)$ is the prior probability.
\end{sym}

\lecture{3}{Independence}

\topic{Event Independence}
\usecase{Use when knowing one event gives no information about another.}
\formulas
\[
A,B\text{ independent}
\quad\Longleftrightarrow\quad
\Prob(A\cap B)=\Prob(A)\Prob(B).
\]
\[
\Prob(A\mid B)=\Prob(A),\qquad \Prob(B)>0.
\]
\begin{sym}
  \item $A,B$ are events.
\end{sym}

\topic{Conditional Independence}
\usecase{Use when independence is checked inside a condition $C$.}
\formulas
\[
\Prob(A\cap B\mid C)=\Prob(A\mid C)\Prob(B\mid C).
\]
\begin{sym}
  \item $C$ is the condition.
\end{sym}

\topic{Mutual Independence}
\usecase{Use when more than two events must be independent together.}
\formulas
\[
\Prob\!\left(\bigcap_{i\in S}A_i\right)
=\prod_{i\in S}\Prob(A_i)
\quad\text{for every subset }S.
\]
\begin{sym}
  \item $S$ is any subset of event indices.
\end{sym}

\lecture{4}{Counting}

\topic{Basic Counting Rule}
\usecase{Use when choices happen in stages.}
\formulas
\[
\#\text{ outcomes}=n_1n_2\cdots n_r.
\]
\begin{sym}
  \item $n_i$ is the number of choices at stage $i$.
  \item $r$ is the number of stages.
\end{sym}

\topic{Permutations}
\usecase{Use when order matters.}
\formulas
\[
P(n,k)=\frac{n!}{(n-k)!}.
\]
\begin{sym}
  \item $n$ is total objects.
  \item $k$ is number of ordered positions.
\end{sym}

\topic{Combinations}
\usecase{Use when order does not matter.}
\formulas
\[
\binom{n}{k}=\frac{n!}{k!(n-k)!}.
\]
\begin{sym}
  \item $n$ is total objects.
  \item $k$ is number selected.
\end{sym}

\topic{Partitions / Multinomial Coefficient}
\usecase{Use when $n$ distinct objects are split into groups of fixed sizes.}
\formulas
\[
\binom{n}{n_1,n_2,\ldots,n_r}
=\frac{n!}{n_1!n_2!\cdots n_r!},
\qquad
\sum_{i=1}^{r}n_i=n.
\]
\begin{sym}
  \item $n_i$ is the size of group $i$.
  \item $r$ is the number of groups.
\end{sym}

\lecture{5}{Random Variables}

\topic{PMF}
\usecase{Use for a discrete random variable.}
\formulas
\[
p_X(x)=\Prob(X=x),\qquad \sum_x p_X(x)=1.
\]
\[
\begin{gathered}
F_X(x)=\Prob(X\le x)=\sum_{u\le x}p_X(u),\\
p_X(k)=F_X(k)-F_X(k-1)\\
\text{for integer-valued }X.
\end{gathered}
\]
\begin{sym}
  \item $X$ is a random variable.
  \item $x$ is a possible value.
  \item $F_X$ is the CDF.
\end{sym}

\topic{Expectation}
\usecase{Use for long-run average value.}
\formulas
\[
\E[X]=\sum_x x\,p_X(x).
\]
\[
\E[g(X)]=\sum_x g(x)p_X(x).
\]
\begin{sym}
  \item $g(X)$ is a transformed random variable.
\end{sym}

\topic{Variance}
\usecase{Use for spread around the mean.}
\formulas
\[
\Var(X)=\E[(X-\E[X])^2]=\E[X^2]-(\E[X])^2.
\]
\[
\Var(aX+b)=a^2\Var(X).
\]
\begin{sym}
  \item $a,b$ are constants.
\end{sym}

\lecture{6}{Conditional PMFs and Total Expectation}

\topic{Conditional PMF}
\usecase{Use for the distribution of $X$ after event $A$.}
\formulas
\[
p_{X\mid A}(x)=
\frac{\Prob(\{X=x\}\cap A)}{\Prob(A)}.
\]
\begin{sym}
  \item $A$ is the event conditioned on.
\end{sym}

\topic{Conditional PMF Given Another RV}
\usecase{Use when $Y=y$ is known.}
\formulas
\[
p_{X\mid Y}(x\mid y)=\frac{p_{X,Y}(x,y)}{p_Y(y)}.
\]
\begin{sym}
  \item $p_{X,Y}$ is the joint PMF.
  \item $p_Y(y)$ is the marginal PMF of $Y$.
\end{sym}

\topic{Total Expectation}
\usecase{Use when expectation is easy after conditioning on cases.}
\formulas
\[
\E[X]=\sum_y \E[X\mid Y=y]\,p_Y(y).
\]
\[
\E[X]=\E[\E[X\mid Y]].
\]
\begin{sym}
  \item $Y$ is the conditioning variable.
\end{sym}

\lecture{7}{Joint PMFs}

\topic{Joint and Marginal PMFs}
\usecase{Use when two discrete variables appear together.}
\formulas
\[
p_{X,Y}(x,y)=\Prob(X=x,Y=y).
\]
\[
p_X(x)=\sum_y p_{X,Y}(x,y),\qquad
p_Y(y)=\sum_x p_{X,Y}(x,y).
\]
\begin{sym}
  \item $p_{X,Y}$ is the joint PMF.
  \item $p_X,p_Y$ are marginal PMFs.
\end{sym}

\topic{Function of Two RVs}
\usecase{Use when $Z=g(X,Y)$.}
\formulas
\[
p_Z(z)=\sum_{(x,y):\,g(x,y)=z}p_{X,Y}(x,y).
\]
\[
\E[g(X,Y)]=\sum_x\sum_y g(x,y)p_{X,Y}(x,y).
\]
\begin{sym}
  \item $Z$ is the derived random variable.
\end{sym}

\topic{Independent Random Variables}
\usecase{Use when joint distribution factors.}
\formulas
\[
p_{X,Y}(x,y)=p_X(x)p_Y(y).
\]
\[
\E[XY]=\E[X]\E[Y],
\qquad
\Var(X+Y)=\Var(X)+\Var(Y).
\]
\begin{sym}
  \item The variance formula needs independence.
\end{sym}

\lecture{8}{Continuous Random Variables}

\topic{PDF and CDF}
\usecase{Use for interval probabilities of continuous random variables.}
\formulas
\[
\Prob(a\le X\le b)=\int_a^b f_X(x)\,dx.
\]
\[
F_X(x)=\Prob(X\le x)=\int_{-\infty}^{x}f_X(t)\,dt,
\qquad
f_X(x)=F_X'(x).
\]
\begin{sym}
  \item $f_X$ is the PDF.
  \item $F_X$ is the CDF.
\end{sym}

\topic{Continuous Expectation and Variance}
\usecase{Use continuous analogs of discrete expectation.}
\formulas
\[
\E[X]=\int_{-\infty}^{\infty}x f_X(x)\,dx,
\qquad
\E[g(X)]=\int_{-\infty}^{\infty}g(x)f_X(x)\,dx.
\]
\[
\Var(X)=\int_{-\infty}^{\infty}(x-\E[X])^2f_X(x)\,dx.
\]
\begin{sym}
  \item Integrals are over the support of $X$.
\end{sym}

\topic{Uniform Distribution}
\usecase{Use when all points in $[a,b]$ are equally dense.}
\formulas
\[
f_X(x)=\frac{1}{b-a},\quad a\le x\le b.
\]
\[
\E[X]=\frac{a+b}{2},\qquad
\Var(X)=\frac{(b-a)^2}{12}.
\]
\begin{sym}
  \item $a,b$ are interval endpoints.
\end{sym}

\topic{Normal Distribution}
\usecase{Use when $X$ is Gaussian or standardized for tables.}
\formulas
\[
X\sim N(\mu,\sigma^2),
\qquad
Z=\frac{X-\mu}{\sigma}\sim N(0,1).
\]
\[
\Prob(X\le x)=\Phi\!\left(\frac{x-\mu}{\sigma}\right).
\]
\begin{sym}
  \item $\mu$ is the mean.
  \item $\sigma^2$ is the variance.
  \item $\Phi$ is the standard normal CDF.
\end{sym}

\topic{Exponential Distribution}
\usecase{Use for continuous memoryless waiting time with rate $\lambda$.}
\formulas
\[
f_X(x)=\lambda e^{-\lambda x},\quad x\ge0,
\qquad
F_X(x)=1-e^{-\lambda x}.
\]
\[
\Prob(X>t)=e^{-\lambda t},
\qquad
\E[X]=\frac{1}{\lambda},
\qquad
\Var(X)=\frac{1}{\lambda^2}.
\]
\[
\Prob(X>s+t\mid X>s)=\Prob(X>t).
\]
\begin{sym}
  \item $\lambda$ is the rate parameter.
  \item $t,x$ are time values.
  \item Last formula is memorylessness.
\end{sym}

\lecture{9}{Joint PDFs and Continuous Conditioning}

\topic{Joint PDF}
\usecase{Use for probabilities over regions in the $(x,y)$ plane.}
\formulas
\[
\Prob((X,Y)\in A)=\iint_A f_{X,Y}(x,y)\,dx\,dy.
\]
\[
f_X(x)=\int_{-\infty}^{\infty}f_{X,Y}(x,y)\,dy,
\quad
f_Y(y)=\int_{-\infty}^{\infty}f_{X,Y}(x,y)\,dx.
\]
\[
\E[g(X,Y)]=\int_{-\infty}^{\infty}\int_{-\infty}^{\infty}
g(x,y)f_{X,Y}(x,y)\,dx\,dy.
\]
\begin{sym}
  \item $A$ is a region in the plane.
  \item $f_X,f_Y$ are marginal PDFs.
\end{sym}

\topic{Conditioning on an Event}
\usecase{Use for a continuous distribution restricted to event $A$.}
\formulas
\[
f_{X\mid A}(x)=
\begin{cases}
\dfrac{f_X(x)}{\Prob(A)}, & x\in A,\\
0, & x\notin A.
\end{cases}
\]
\[
\Prob(A)=\int_A f_X(x)\,dx.
\]
\begin{sym}
  \item $A$ is the conditioning event.
\end{sym}

\topic{Continuous Total Probability}
\usecase{Use when cases $A_i$ partition the sample space.}
\formulas
\[
f_X(x)=\sum_i \Prob(A_i)f_{X\mid A_i}(x).
\]
\[
f_X(x)=\int f_Y(y)f_{X\mid Y}(x\mid y)\,dy.
\]
\begin{sym}
  \item $A_i$ are disjoint exhaustive cases.
  \item $Y$ is a continuous conditioning variable.
\end{sym}

\topic{Conditional PDF}
\usecase{Use when one continuous variable is observed.}
\formulas
\[
f_{X\mid Y}(x\mid y)=\frac{f_{X,Y}(x,y)}{f_Y(y)}.
\]
\[
\E[X\mid Y=y]=\int x f_{X\mid Y}(x\mid y)\,dx.
\]
\begin{sym}
  \item $f_Y(y)>0$ is required.
\end{sym}

\topic{Continuous Independence}
\usecase{Use when joint PDF factors into marginals.}
\formulas
\[
f_{X,Y}(x,y)=f_X(x)f_Y(y).
\]
\begin{sym}
  \item The support must also factor consistently.
\end{sym}

\lecture{10}{Derived Distributions and Continuous Bayes}

\topic{Monotone Transformation}
\usecase{Use when $Y=g(X)$ and $g$ is one-to-one.}
\formulas
\[
f_Y(y)=f_X(g^{-1}(y))
\left|\frac{d}{dy}g^{-1}(y)\right|.
\]
\[
Y=aX+b
\quad\Longrightarrow\quad
f_Y(y)=\frac{1}{|a|}
f_X\!\left(\frac{y-b}{a}\right).
\]
\begin{sym}
  \item $g^{-1}(y)$ is the inverse value of $x$.
  \item $a,b$ are constants and $a\ne0$.
\end{sym}

\topic{CDF Method}
\usecase{Use when transformation is not one-to-one.}
\formulas
\[
F_Y(y)=\Prob(Y\le y)=\Prob(g(X)\le y),
\qquad
f_Y(y)=\frac{d}{dy}F_Y(y).
\]
\begin{sym}
  \item $Y=g(X)$ is the transformed variable.
\end{sym}

\topic{Max and Min of Independent RVs}
\usecase{Use when a derived variable is the maximum or minimum.}
\formulas
\[
M=\max(X_1,\ldots,X_n)
\quad\Longrightarrow\quad
F_M(m)=\prod_{i=1}^{n}F_{X_i}(m).
\]
\[
L=\min(X_1,\ldots,X_n)
\quad\Longrightarrow\quad
\Prob(L>\ell)=\prod_{i=1}^{n}\Prob(X_i>\ell).
\]
\[
F_L(\ell)=1-\prod_{i=1}^{n}(1-F_{X_i}(\ell)).
\]
\begin{sym}
  \item $M$ is the maximum.
  \item $L$ is the minimum.
  \item Independence is required.
\end{sym}

\topic{Continuous Bayes}
\usecase{Use when evidence is continuous.}
\formulas
\[
f_{X\mid Y}(x\mid y)=
\frac{f_{Y\mid X}(y\mid x)f_X(x)}
{\int f_{Y\mid X}(y\mid u)f_X(u)\,du}.
\]
\[
p_{X\mid Y}(x\mid y)=
\frac{p_X(x)f_{Y\mid X}(y\mid x)}
{\sum_u p_X(u)f_{Y\mid X}(y\mid u)}.
\]
\[
\Prob(A\mid Y=y)=
\frac{\Prob(A)f_{Y\mid A}(y)}
{\Prob(A)f_{Y\mid A}(y)+\Prob(A^c)f_{Y\mid A^c}(y)}.
\]
\begin{sym}
  \item First formula is continuous hidden $X$.
  \item Second formula is discrete hidden $X$.
  \item Last formula is Bayes for an event and continuous evidence.
\end{sym}

\lecture{11}{Convolution, Covariance, Correlation}

\topic{Convolution}
\usecase{Use for sum of independent random variables.}
\formulas
\[
W=X+Y,\quad X\perp Y.
\]
\[
p_W(w)=\sum_x p_X(x)p_Y(w-x).
\]
\[
f_W(w)=\int_{-\infty}^{\infty}f_X(x)f_Y(w-x)\,dx.
\]
\begin{sym}
  \item $w-x$ is the value needed from $Y$.
\end{sym}

\topic{Mixed Discrete--Continuous Sum}
\usecase{Use when $X$ is discrete and $Y$ is continuous.}
\formulas
\[
f_W(w)=\sum_x p_X(x)f_Y(w-x).
\]
\begin{sym}
  \item $W=X+Y$.
\end{sym}

\topic{Covariance and Correlation}
\usecase{Use for linear co-movement between variables.}
\formulas
\[
\Cov(X,Y)=\E[XY]-\E[X]\E[Y].
\]
\[
\Corr(X,Y)=\frac{\Cov(X,Y)}{\sigma_X\sigma_Y}.
\]
\[
\Var(X+Y)=\Var(X)+\Var(Y)+2\Cov(X,Y).
\]
\begin{sym}
  \item $\sigma_X=\sqrt{\Var(X)}$.
\end{sym}

\lecture{12}{Conditional Expectation and Total Variance}

\topic{Conditional Expectation}
\usecase{Use for the mean after observing $Y=y$.}
\formulas
\[
\E[X\mid Y=y]=\sum_x x p_{X\mid Y}(x\mid y).
\]
\[
\E[X\mid Y=y]=\int x f_{X\mid Y}(x\mid y)\,dx.
\]
\begin{sym}
  \item Use the sum for discrete $X$.
  \item Use the integral for continuous $X$.
\end{sym}

\topic{Total Variance}
\usecase{Use when variance splits into within-case and between-case parts.}
\formulas
\[
\Var(X)=\E[\Var(X\mid Y)]+\Var(\E[X\mid Y]).
\]
\begin{sym}
  \item $\E[\Var(X\mid Y)]$ is average within-case variance.
  \item $\Var(\E[X\mid Y])$ is variance of case means.
\end{sym}

\topic{Random Sum}
\usecase{Use when adding a random number of i.i.d. terms.}
\formulas
\[
S_N=X_1+\cdots+X_N,
\qquad N\perp \{X_i\}.
\]
\[
\E[S_N]=\E[N]\E[X],
\]
\[
\Var(S_N)=\E[N]\Var(X)+\Var(N)(\E[X])^2.
\]
\begin{sym}
  \item $N$ is the random number of terms.
  \item $X_i$ are i.i.d. copies of $X$.
\end{sym}

\lecture{13}{Bernoulli Process}

\topic{Bernoulli and Binomial}
\usecase{Use Binomial when the total number of trials is fixed.}
\formulas
\[
X_i\sim \operatorname{Bernoulli}(p),
\qquad
\Prob(X_i=1)=p.
\]
\[
S_n=\sum_{i=1}^{n}X_i\sim \operatorname{Binomial}(n,p),
\]
\[
\Prob(S_n=k)=\binom{n}{k}p^k(1-p)^{n-k}.
\]
\[
\E[S_n]=np,\qquad \Var(S_n)=np(1-p).
\]
\begin{sym}
  \item $n$ is fixed number of trials.
  \item $k$ is number of successes.
  \item $p$ is success probability per trial.
\end{sym}

\topic{First Success}
\usecase{Use Geometric for the trial number of the first success.}
\formulas
\[
T_1\sim \operatorname{Geometric}(p),
\qquad
\Prob(T_1=t)=(1-p)^{t-1}p.
\]
\[
\E[T_1]=\frac{1}{p},
\qquad
\Var(T_1)=\frac{1-p}{p^2}.
\]
\[
\Prob(T_1>s+t\mid T_1>s)=\Prob(T_1>t).
\]
\begin{sym}
  \item $T_1$ counts the success trial.
  \item Last formula is memorylessness.
\end{sym}

\topic{Kth Success}
\usecase{Use Negative Binomial/Pascal for the trial of the kth success.}
\formulas
\[
Y_k\sim \operatorname{NegativeBinomial}(k,p),
\]
\[
\Prob(Y_k=t)=\binom{t-1}{k-1}p^k(1-p)^{t-k},
\quad t\ge k.
\]
\[
\E[Y_k]=\frac{k}{p},
\qquad
\Var(Y_k)=\frac{k(1-p)}{p^2}.
\]
\begin{sym}
  \item $Y_k$ is the trial of the kth success.
  \item The final trial $t$ is a success.
\end{sym}

\topic{Splitting and Merging Bernoulli Processes}
\usecase{Use for per-slot arrivals that are filtered or combined.}
\formulas
\[
\Prob(\text{type }i\text{ success in a slot})=p\,q_i.
\]
\[
\Prob(\text{merged arrival})
=1-\prod_i(1-p_i).
\]
\[
\text{For two streams:}\quad
1-(1-p)(1-q)=p+q-pq.
\]
\begin{sym}
  \item $p$ is original success probability.
  \item $q_i$ is type/filter probability.
  \item $p_i$ are independent per-slot arrival probabilities.
\end{sym}

\lecture{14}{Poisson Process I}

\topic{Poisson Counts}
\usecase{Use for continuous-time arrivals with rate $\lambda$.}
\formulas
\[
N(t)\sim \operatorname{Poisson}(\lambda t),
\qquad
\Prob(N(t)=k)=e^{-\lambda t}\frac{(\lambda t)^k}{k!}.
\]
\[
\E[N(t)]=\lambda t,\qquad \Var(N(t))=\lambda t.
\]
\begin{sym}
  \item $N(t)$ is number of arrivals in length $t$.
  \item $\lambda$ is arrival rate per unit time.
\end{sym}

\topic{Small Interval Approximation}
\usecase{Use for local behavior of a Poisson process.}
\formulas
\[
\Prob(N(\delta)=0)\approx 1-\lambda\delta,
\quad
\Prob(N(\delta)=1)\approx \lambda\delta,
\quad
\Prob(N(\delta)\ge 2)\approx 0.
\]
\begin{sym}
  \item $\delta$ is a very small time length.
\end{sym}

\topic{Independent Increments}
\usecase{Use for counts in non-overlapping intervals.}
\formulas
\[
N(t_2)-N(t_1)\sim
\operatorname{Poisson}(\lambda(t_2-t_1)).
\]
\[
N(t_2)-N(t_1)
\text{ is independent of disjoint interval counts.}
\]
\begin{sym}
  \item $t_2-t_1$ is the interval length.
\end{sym}

\lecture{15}{Poisson Process II}

\topic{First Arrival Time}
\usecase{Use Exponential for waiting time to the first Poisson arrival.}
\formulas
\[
T_1\sim \operatorname{Exponential}(\lambda),
\qquad
f_{T_1}(t)=\lambda e^{-\lambda t},\quad t\ge 0.
\]
\[
\Prob(T_1>t)=e^{-\lambda t},
\qquad
\E[T_1]=\frac{1}{\lambda},
\qquad
\Var(T_1)=\frac{1}{\lambda^2}.
\]
\begin{sym}
  \item $T_1$ is time until first arrival.
\end{sym}

\topic{Kth Arrival Time}
\usecase{Use Erlang for waiting time to the kth Poisson arrival.}
\formulas
\[
Y_k\sim \operatorname{Erlang}(k,\lambda),
\qquad
f_{Y_k}(y)=\frac{\lambda^k y^{k-1}e^{-\lambda y}}{(k-1)!}.
\]
\[
\E[Y_k]=\frac{k}{\lambda},
\qquad
\Var(Y_k)=\frac{k}{\lambda^2}.
\]
\begin{sym}
  \item $Y_k$ is time of the kth arrival.
\end{sym}

\topic{Arrival Time and Count Conversion}
\usecase{Use when converting kth-arrival questions into count questions.}
\formulas
\[
\Prob(Y_k\le t)=\Prob(N(t)\ge k),
\qquad
\Prob(Y_k>t)=\Prob(N(t)\le k-1).
\]
\begin{sym}
  \item $N(t)$ is the count by time $t$.
\end{sym}

\topic{Merging and Splitting Poisson Processes}
\usecase{Use merging for combined streams and splitting for typed arrivals.}
\formulas
\[
\begin{gathered}
N_i(t)\sim\operatorname{Poisson}(\lambda_i t)
\text{ independent},\\
\sum_i N_i(t)\sim
\operatorname{Poisson}\!\left(\sum_i\lambda_i\,t\right).
\end{gathered}
\]
\[
\Prob(\text{next arrival from stream }i)=
\frac{\lambda_i}{\sum_j\lambda_j}.
\]
\[
N_i(t)\sim \operatorname{Poisson}(\lambda p_i t)
\quad\text{after splitting by type }i.
\]
\begin{sym}
  \item $\lambda_i$ is rate of stream $i$.
  \item $p_i$ is type probability after splitting.
\end{sym}

\topic{Conditioning on Total Poisson Count}
\usecase{Use Binomial-style counting once total arrivals are fixed.}
\formulas
\[
N_1(t)\mid N(t)=m\sim \operatorname{Binomial}(m,p_1).
\]
\begin{sym}
  \item $m$ is the given total count.
  \item $p_1$ is probability of type 1.
\end{sym}

\topic{Random Incidence}
\usecase{Use when observing a renewal process at a random time.}
\formulas
\[
\Prob(\text{chosen interval has length }\ell_i)
=\frac{\ell_i p_i}{\E[L]}.
\]
\[
\E[\text{forward wait}\mid L=\ell_i]=\frac{\ell_i}{2}.
\]
\begin{sym}
  \item $L$ is ordinary interarrival length.
  \item $p_i=\Prob(L=\ell_i)$.
\end{sym}

\topic{Last of Several Exponential Lifetimes}
\usecase{Use when waiting until all $n$ independent exponential lifetimes have ended.}
\formulas
\[
\min(T_1,\ldots,T_n)\sim \operatorname{Exponential}(n\lambda)
\quad\text{if }T_i\sim\operatorname{Exponential}(\lambda).
\]
\[
\E[\max(T_1,\ldots,T_n)]
=\frac{1}{\lambda}\sum_{j=1}^{n}\frac{1}{j}.
\]
\begin{sym}
  \item $T_i$ are independent exponential lifetimes.
  \item $\lambda$ is the lifetime rate.
\end{sym}

\lecture{16}{Markov Processes I}

\topic{Markov Property}
\usecase{Use when the next state depends only on the current state.}
\formulas
\[
\Prob(X_{n+1}=j\mid X_n=i,X_{n-1},\ldots,X_0)
=p_{ij}.
\]
\[
p_{ij}=\Prob(X_{n+1}=j\mid X_n=i).
\]
\begin{sym}
  \item $X_n$ is the state at time $n$.
  \item $p_{ij}$ is one-step transition probability.
\end{sym}

\topic{N-Step Transitions}
\usecase{Use for probability of being in state $j$ after $n$ steps.}
\formulas
\[
r_{ij}(n)=\Prob(X_n=j\mid X_0=i)=(P^n)_{ij}.
\]
\[
r_{ij}(n)=\sum_k r_{ik}(n-1)p_{kj}.
\]
\begin{sym}
  \item $P$ is the transition matrix.
  \item $r_{ij}(n)$ is the $n$-step transition probability.
\end{sym}

\topic{State Occupancy}
\usecase{Use when the initial state is random.}
\formulas
\[
\pi^{(n)}=\pi^{(0)}P^n,
\qquad
\Prob(X_n=j)=\sum_i \pi_i^{(0)}r_{ij}(n).
\]
\begin{sym}
  \item $\pi^{(n)}$ is the row vector of state probabilities at time $n$.
\end{sym}

\lecture{17}{Markov Processes II}

\topic{Steady-State Probabilities}
\usecase{Use for long-run state probabilities in an aperiodic single recurrent class.}
\formulas
\[
\lim_{n\to\infty}r_{ij}(n)=\pi_j.
\]
\[
\pi_j=\sum_i \pi_i p_{ij},\qquad \sum_j \pi_j=1.
\]
\[
\pi=\pi P.
\]
\begin{sym}
  \item $\pi_j$ is steady-state probability of state $j$.
\end{sym}

\topic{Long-Run Transition Frequency}
\usecase{Use for fraction of transitions from $i$ to $j$.}
\formulas
\[
\text{long-run frequency of }i\to j=\pi_i p_{ij}.
\]
\begin{sym}
  \item $\pi_i$ is long-run fraction of time in state $i$.
\end{sym}

\topic{Birth--Death Chains}
\usecase{Use for nearest-neighbor Markov chains.}
\formulas
\[
\pi_i p_i=\pi_{i+1}q_{i+1}.
\]
\[
\pi_i=\rho^i\pi_0,\qquad \rho=\frac{p}{q}\quad\text{when }p_i=p,\ q_i=q.
\]
\[
\pi_0=\frac{1-\rho}{1-\rho^{m+1}}\quad(\rho\ne1,\ 0\le i\le m).
\]
\[
\pi_i=(1-\rho)\rho^i\quad(\rho<1,\ i=0,1,\ldots).
\]
\begin{sym}
  \item $p_i$ is transition probability from $i$ to $i+1$.
  \item $q_i$ is transition probability from $i$ to $i-1$.
  \item $\rho$ is the load ratio.
\end{sym}

\lecture{18}{Markov Processes III}

\topic{Phone Blocking Model}
\usecase{Use for $B$ lines, Poisson arrivals, exponential call durations.}
\formulas
\[
\lambda\pi_{i-1}=i\mu\pi_i,\qquad i=1,\ldots,B.
\]
\[
\pi_i=\frac{(\lambda/\mu)^i}{i!}\pi_0,
\qquad
\pi_0=\left(\sum_{i=0}^{B}\frac{(\lambda/\mu)^i}{i!}\right)^{-1}.
\]
\[
\Prob(\text{blocked})=\pi_B.
\]
\begin{sym}
  \item $\lambda$ is call arrival rate.
  \item $\mu$ is service completion rate per active call.
  \item $B$ is number of lines.
\end{sym}

\topic{Absorption Probability}
\usecase{Use for probability of eventually reaching an absorbing target.}
\formulas
\[
a_i=\sum_j p_{ij}a_j
\quad\text{for nonabsorbing states }i.
\]
\[
\begin{gathered}
a_i=1\quad\text{on target absorbing states},\\
a_i=0\quad\text{on other absorbing states}.
\end{gathered}
\]
\begin{sym}
  \item $a_i$ is absorption probability starting from state $i$.
\end{sym}

\topic{Expected Time to Absorption}
\usecase{Use for expected number of transitions until absorption.}
\formulas
\[
\begin{gathered}
t_i=1+\sum_j p_{ij}t_j
\quad\text{for nonabsorbing states }i,\\
t_i=0\quad\text{on absorbing states}.
\end{gathered}
\]
\begin{sym}
  \item $t_i$ is expected time to absorption from state $i$.
\end{sym}

\topic{Mean First Passage and Recurrence}
\usecase{Use for expected time to hit a state $s$.}
\formulas
\[
t_s=0,\qquad
t_i=1+\sum_j p_{ij}t_j,\quad i\ne s.
\]
\[
t_s^*=1+\sum_j p_{sj}t_j.
\]
\[
t_s^*=\frac{1}{\pi_s}
\quad\text{in steady state for a positive recurrent state.}
\]
\begin{sym}
  \item $t_i$ is mean first passage time from $i$ to $s$.
  \item $t_s^*$ is mean recurrence time of $s$.
  \item $\pi_s$ is the steady-state probability of $s$.
\end{sym}

\lecture{19}{Chebyshev and WLLN}

\topic{Chebyshev's Inequality}
\usecase{Use for probability bounds using only mean and variance.}
\formulas
\[
\Prob(|X-\mu|\ge c)\le \frac{\sigma^2}{c^2}.
\]
\[
\Prob(|X-\mu|\ge k\sigma)\le \frac{1}{k^2}.
\]
\begin{sym}
  \item $\mu=\E[X]$.
  \item $\sigma^2=\Var(X)$.
\end{sym}

\topic{Convergence in Probability}
\usecase{Use when random variables concentrate near a number.}
\formulas
\[
Y_n\xrightarrow{p}a
\quad\Longleftrightarrow\quad
\lim_{n\to\infty}\Prob(|Y_n-a|\ge \epsilon)=0
\quad\forall \epsilon>0.
\]
\begin{sym}
  \item $a$ is the deterministic limit.
\end{sym}

\topic{Weak Law of Large Numbers}
\usecase{Use for the sample mean of i.i.d. variables.}
\formulas
\[
M_n=\frac{X_1+\cdots+X_n}{n},
\qquad
\E[M_n]=\mu,\qquad
\Var(M_n)=\frac{\sigma^2}{n}.
\]
\[
\Prob(|M_n-\mu|\ge \epsilon)
\le \frac{\sigma^2}{n\epsilon^2},
\qquad
M_n\xrightarrow{p}\mu.
\]
\begin{sym}
  \item $X_i$ are i.i.d. with mean $\mu$ and variance $\sigma^2$.
\end{sym}

\lecture{20}{Central Limit Theorem}

\topic{CLT Standardization}
\usecase{Use for sums of many i.i.d. variables with finite variance.}
\formulas
\[
S_n=X_1+\cdots+X_n,
\qquad
Z_n=\frac{S_n-n\mu}{\sigma\sqrt{n}}.
\]
\[
\Prob(Z_n\le z)\to \Phi(z).
\]
\[
S_n\approx N(n\mu,n\sigma^2),
\qquad
M_n\approx N\!\left(\mu,\frac{\sigma^2}{n}\right).
\]
\begin{sym}
  \item $\Phi$ is the standard normal CDF.
  \item $M_n=S_n/n$.
\end{sym}

\topic{Normal Approximation to Binomial}
\usecase{Use for large $n$ Binomial counts.}
\formulas
\[
S_n\sim \operatorname{Binomial}(n,p),
\qquad
\frac{S_n-np}{\sqrt{np(1-p)}}\approx N(0,1).
\]
\[
\Prob(S_n\le m)\approx
\Phi\!\left(\frac{m+0.5-np}{\sqrt{np(1-p)}}\right).
\]
\begin{sym}
  \item $0.5$ is the continuity correction.
\end{sym}

\topic{Poisson Approximation and Normal Approximation}
\usecase{Use Poisson for rare Binomial events and Normal for large counts.}
\formulas
\[
\operatorname{Binomial}(n,p)\approx \operatorname{Poisson}(\lambda),
\qquad
\lambda=np
\quad(n\text{ large},\ p\text{ small}).
\]
\[
\operatorname{Poisson}(\lambda)\approx N(\lambda,\lambda)
\quad(\lambda\text{ large}).
\]
\begin{sym}
  \item $\lambda$ is the Poisson mean and variance.
\end{sym}

\lecture{21}{Bayesian Inference}

\topic{Bayesian Posterior}
\usecase{Use when an unknown parameter is modeled as random.}
\formulas
\[
p_{\Theta\mid X}(\theta\mid x)=
\frac{p_\Theta(\theta)p_{X\mid\Theta}(x\mid\theta)}
{\sum_u p_\Theta(u)p_{X\mid\Theta}(x\mid u)}.
\]
\[
f_{\Theta\mid X}(\theta\mid x)=
\frac{f_\Theta(\theta)f_{X\mid\Theta}(x\mid\theta)}
{\int f_\Theta(u)f_{X\mid\Theta}(x\mid u)\,du}.
\]
\begin{sym}
  \item $\Theta$ is the unknown random parameter.
  \item $X=x$ is the observed data.
\end{sym}

\topic{MAP and LMS Estimates}
\usecase{Use MAP for most likely posterior value and LMS for squared-error loss.}
\formulas
\[
\hat{\theta}_{MAP}=\arg\max_\theta p_{\Theta\mid X}(\theta\mid x)
\quad\text{or}\quad
\arg\max_\theta f_{\Theta\mid X}(\theta\mid x).
\]
\[
\hat{\Theta}_{LMS}=\E[\Theta\mid X].
\]
\begin{sym}
  \item MAP maximizes posterior probability/density.
  \item LMS minimizes mean squared error.
\end{sym}

\lecture{22}{LMS and Linear LMS Estimation}

\topic{Conditional Mean as LMS}
\usecase{Use for the best estimator under mean squared error.}
\formulas
\[
\hat{\Theta}=\E[\Theta\mid X]
=\arg\min_{g(X)}\E[(\Theta-g(X))^2].
\]
\[
\E[(\Theta-\E[\Theta\mid X])^2\mid X=x]
=\Var(\Theta\mid X=x).
\]
\begin{sym}
  \item $g(X)$ is any estimator based on $X$.
\end{sym}

\topic{LMS Error Properties}
\usecase{Use for orthogonality and variance decomposition.}
\formulas
\[
\tilde{\Theta}=\Theta-\hat{\Theta},
\qquad
\E[\tilde{\Theta}]=0,
\qquad
\E[\tilde{\Theta}h(X)]=0.
\]
\[
\Var(\Theta)=\Var(\hat{\Theta})+\Var(\tilde{\Theta}).
\]
\begin{sym}
  \item $h(X)$ is any function of the observation.
\end{sym}

\topic{Linear LMS}
\usecase{Use when restricting the estimator to be linear in $X$.}
\formulas
\[
\hat{\Theta}_L=\E[\Theta]
\;+\;\frac{\Cov(\Theta,X)}{\Var(X)}(X-\E[X]).
\]
\[
\hat{\Theta}_L=\E[\Theta]
\;+\;\Cov(\Theta,\mathbf{X})\Var(\mathbf{X})^{-1}
(\mathbf{X}-\E[\mathbf{X}]).
\]
\begin{sym}
  \item First formula is one observation.
  \item Second formula is vector observations.
\end{sym}

\lecture{23}{Classical Statistics}

\topic{Maximum Likelihood}
\usecase{Use when the parameter is unknown but not random.}
\formulas
\[
L(\theta)=p_X(x;\theta)
\quad\text{or}\quad
L(\theta)=f_X(x;\theta).
\]
\[
\hat{\theta}_{ML}=\arg\max_\theta L(\theta)
=\arg\max_\theta \log L(\theta).
\]
\[
L(\theta)=\prod_{i=1}^{n}p_X(x_i;\theta)
\quad\text{or}\quad
\prod_{i=1}^{n}f_X(x_i;\theta).
\]
\begin{sym}
  \item $\theta$ is a fixed unknown parameter.
  \item $x_i$ are observed data.
\end{sym}

\topic{Estimator Properties}
\usecase{Use to judge bias, consistency, and mean squared error.}
\formulas
\[
\text{Unbiased:}\quad \E_\theta[\hat{\theta}_n]=\theta.
\]
\[
\text{Consistent:}\quad \hat{\theta}_n\xrightarrow{p}\theta.
\]
\[
\operatorname{MSE}(\hat{\theta})
=\E[(\hat{\theta}-\theta)^2]
=\Var(\hat{\theta})+\operatorname{bias}(\hat{\theta})^2.
\]
\begin{sym}
  \item $\operatorname{bias}(\hat{\theta})=\E[\hat{\theta}]-\theta$.
\end{sym}

\topic{Sample Mean and Confidence Interval}
\usecase{Use for estimating an unknown mean.}
\formulas
\[
\bar X=\frac{X_1+\cdots+X_n}{n},
\qquad
\E[\bar X]=\mu,\qquad
\Var(\bar X)=\frac{\sigma^2}{n}.
\]
\[
\bar X\pm z_{\alpha/2}\frac{\sigma}{\sqrt n}
\quad\text{is a }(1-\alpha)\text{ CI when }\sigma\text{ is known.}
\]
\[
\bar X\pm z_{\alpha/2}\frac{S_n}{\sqrt n}
\quad\text{large-sample plug-in CI.}
\]
\[
S_n^2=\frac{1}{n-1}\sum_{i=1}^{n}(X_i-\bar X)^2.
\]
\begin{sym}
  \item $z_{\alpha/2}$ satisfies $\Prob(Z>z_{\alpha/2})=\alpha/2$.
  \item $S_n^2$ estimates $\sigma^2$.
\end{sym}

\lecture{24}{Regression and Likelihood Ratio Tests}

\topic{Simple Linear Regression}
\usecase{Use to fit a line $y=\beta_0+\beta_1 x$ by least squares.}
\formulas
\[
(\hat{\beta}_0,\hat{\beta}_1)
=\arg\min_{\beta_0,\beta_1}
\sum_{i=1}^{n}(y_i-\beta_0-\beta_1x_i)^2.
\]
\[
\hat{\beta}_1=
\frac{\sum_i (x_i-\bar x)(y_i-\bar y)}
{\sum_i (x_i-\bar x)^2},
\qquad
\hat{\beta}_0=\bar y-\hat{\beta}_1\bar x.
\]
\begin{sym}
  \item $(x_i,y_i)$ are data points.
  \item $\hat y_i=\hat\beta_0+\hat\beta_1x_i$.
\end{sym}

\topic{Likelihood Ratio Test}
\usecase{Use for binary hypothesis testing between $H_0$ and $H_1$.}
\formulas
\[
\Lambda(x)=\frac{p_X(x\mid H_1)}{p_X(x\mid H_0)}
\quad\text{or}\quad
\frac{f_X(x\mid H_1)}{f_X(x\mid H_0)}.
\]
\[
\text{Choose }H_1\text{ if }\Lambda(x)>\eta.
\]
\[
\text{MAP threshold:}\quad
\eta=\frac{\Prob(H_0)}{\Prob(H_1)}.
\]
\begin{sym}
  \item $\eta$ controls the tradeoff between error types.
\end{sym}

\topic{Error Probabilities}
\usecase{Use to measure false alarm and missed detection.}
\formulas
\[
\alpha=\Prob(\text{reject }H_0\mid H_0),
\qquad
\beta=\Prob(\text{accept }H_0\mid H_1).
\]
\begin{sym}
  \item $\alpha$ is Type I error.
  \item $\beta$ is Type II error.
\end{sym}

\lecture{25}{Hypothesis Testing}

\topic{Normal Mean Test}
\usecase{Use LRT for normal data with different means and known variance.}
\formulas
\[
H_0:X_i\sim N(\mu_0,\sigma^2),
\qquad
H_1:X_i\sim N(\mu_1,\sigma^2).
\]
\[
\log\Lambda(x)
=\frac{\mu_1-\mu_0}{\sigma^2}\sum_{i=1}^n x_i
-\frac{n(\mu_1^2-\mu_0^2)}{2\sigma^2}.
\]
\[
\mu_1>\mu_0:\quad \text{reject }H_0\text{ if }\sum_i X_i>c.
\]
\begin{sym}
  \item $c$ is chosen from the desired significance level.
\end{sym}

\topic{Normal Variance Test}
\usecase{Use when testing variance of zero-mean normal data.}
\formulas
\[
H_0:X_i\sim N(0,\sigma_0^2),
\qquad
H_1:X_i\sim N(0,\sigma_1^2).
\]
\[
\sum_{i=1}^{n}\frac{X_i^2}{\sigma_0^2}\sim \chi_n^2
\quad\text{under }H_0.
\]
\begin{sym}
  \item $\chi_n^2$ is chi-square with $n$ degrees of freedom.
\end{sym}

\topic{Chi-Square Goodness of Fit}
\usecase{Use to test whether observed bin counts match expected probabilities.}
\formulas
\[
T=\sum_{i=1}^{k}\frac{(N_i-np_i)^2}{np_i}.
\]
\[
T\approx \chi_{k-1}^2
\quad\text{under }H_0\text{ when }n\text{ is large.}
\]
\begin{sym}
  \item $N_i$ is observed count in bin $i$.
  \item $np_i$ is expected count in bin $i$.
  \item $k$ is number of bins.
\end{sym}

\topic{Kolmogorov--Smirnov Statistic}
\usecase{Use to compare an empirical CDF with a proposed CDF.}
\formulas
\[
D_n=\max_x |F_X(x)-\hat F_X(x)|.
\]
\begin{sym}
  \item $F_X$ is the proposed CDF.
  \item $\hat F_X$ is the empirical CDF.
\end{sym}

\section*{Quick Symbol Reference}
\begin{tabular}{p{0.55in}p{2.65in}}
$n$ & number of trials, observations, or steps\\
$k$ & successes, arrivals, bins, or target count\\
$p$ & Bernoulli success probability\\
$\lambda$ & Poisson arrival rate\\
$\mu$ & mean, or service rate in queueing context\\
$\sigma^2$ & variance\\
$N(t)$ & Poisson count by time $t$\\
$P$ & Markov transition matrix\\
$\pi$ & steady-state distribution\\
$\theta$ & unknown parameter\\
$\hat\theta$ & estimator of $\theta$\\
$\alpha$ & Type I error / significance level\\
$\beta$ & Type II error\\
\end{tabular}

\end{document}
